% =============================================================
%  Section 4 — Five-Signal Detection Framework
% =============================================================

\section{Five-Signal Detection Framework}\label{sec:detection}

The threat model of \cref{sec:threat} identified eight attack chains
that exploit the gap between human behavioral assumptions and agent
capabilities.  We now present a detection framework built on a
fundamentally different premise: rather than modeling what humans
\emph{do} and flagging deviations, we model what agents
\emph{cannot avoid doing} and detect the presence of those invariant
signatures.


% ============================================================
\subsection{Design Principles}\label{subsec:principles}

The framework rests on a paradigm shift from \emph{human-behavioral}
detection to \emph{agent-invariant} detection.  The operative question
changes from ``does this look like a human?'' to ``does this exhibit
properties that are \emph{present} in agents but \emph{absent} in
humans?''

This shift is not cosmetic.  Human-behavioral models degrade as agents
learn to mimic human patterns (CHAIN\_6 in \cref{sec:threat}), whereas
agent-invariant signals exploit properties that agents cannot suppress
without sacrificing the economic advantage that motivates the attack in
the first place.  We formalize this intuition through four necessity
principles:

\begin{enumerate}[label=\textbf{P\arabic*.},leftmargin=2.5em]
  \item \textbf{Economic Necessity.}\label{principle:econ}
    An agent must transact at volumes or values that justify its
    operational cost.  Transactions that are economically irrational
    under any plausible human utility function---but rational under
    agent amortization---constitute a detectable signal.

  \item \textbf{Network Necessity.}\label{principle:net}
    An agent must interact with counterparties to accomplish its
    objective, producing graph structures (hubs, fans, chains) that
    differ from organic human social graphs.

  \item \textbf{Temporal Necessity.}\label{principle:temp}
    An agent operates on machine time.  Even when throttled to
    mimic human cadence, the residual timing distribution reveals
    either machine-speed precision ($<400$\,ms intervals) or
    machine-generated regularity ($\cvt < 0.05$)---both absent in
    genuine human or legitimate-agent traffic.

  \item \textbf{Value Flow Necessity.}\label{principle:vf}
    Illicit value must move from source to destination.
    Regardless of the layering strategy, the aggregate flow velocity
    and directionality leave detectable traces in the value-flow
    graph.
\end{enumerate}

\noindent
Each principle maps to one or more of the five signals described
below.  A fifth signal, Cross-Platform Correlation, anticipates the
multi-chain setting of CHAIN\_5 but is not active in the current
single-chain implementation.


% ============================================================
\subsection{Agent-Invariant Signals}\label{subsec:signals}

We define five composite signals, each decomposed into sub-signals
with empirically derived weights.
\Cref{tab:signal_overview} provides a summary; the sub-signal
specifications and scoring functions are detailed in
\cref{app:signals}.

\begin{table}[t]
  \centering
  \caption{%
    Five-signal overview.  Sub-signal weights were derived from
    real on-chain validation data (81,904 Base chain USDC transactions;
    see~\cref{sec:evaluation}).  CP is architecturally defined but
    inactive in the current single-chain deployment.%
  }
  \label{tab:signal_overview}
  \small
  \begin{tabularx}{\textwidth}{@{}l l X l@{}}
    \toprule
    \textbf{Signal} & \textbf{Abbr.} & \textbf{Sub-signals (weight)}
      & \textbf{Principle} \\
    \midrule
    Economic Rationality & ER
      & Utility maximization / circular flow (40\%),
        purpose coherence (35\%),
        concentration patterns (25\%)
      & P1 \\[4pt]
    Network Topology & NT
      & Sender centrality (50\%),
        receiver centrality (30\%),
        path anomaly (20\%)
      & P2 \\[4pt]
    Value Flow & VF
      & Flow velocity (50\%),
        layering indicator (50\%)
      & P4 \\[4pt]
    Temporal Consistency & TC
      & Circadian violation (35\%),
        timing regularity (35\%),
        burst detection (30\%)
      & P3 \\[4pt]
    Cross-Platform Correlation & CP
      & \emph{Inactive}---designed for simultaneous
        multi-chain presence detection
      & P2, P3 \\
    \bottomrule
  \end{tabularx}
\end{table}


% ---- Signal 1: Economic Rationality ----
\subsubsection{Signal 1: Economic Rationality (ER)}\label{subsub:er}

The Economic Rationality signal detects transactions that are
economically incoherent under human decision-making but perfectly
rational for an amortized autonomous agent.

\paragraph{Sub-signals.}
\emph{Utility maximization} flags circular flows---value that
departs from and returns to the same address through intermediate
hops---which serve no legitimate economic purpose but are a hallmark
of wash trading and layering.
\emph{Purpose coherence} measures whether an address's transaction
mix is consistent with any identifiable economic role (merchant,
consumer, liquidity provider); addresses with incoherent mixes score
higher.
\emph{Concentration patterns} detect addresses whose counterparty
set is anomalously narrow or anomalously broad relative to their
transaction volume.

\paragraph{Agent-economic threshold.}
A key empirical finding is that agents routinely transact at values
between \$0.06 and \$0.91---amounts that fall below the cognitive
cost threshold at which a human would bother initiating a
transaction.  The marginal cost of human attention to a financial
decision has been estimated at \$1--\$5~\cite{kahneman2011};
agent amortization reduces this to effectively zero.  Transactions
below this threshold therefore carry a strong \emph{a~priori}
probability of agent origin, independent of any behavioral feature.


% ---- Signal 2: Network Topology ----
\subsubsection{Signal 2: Network Topology (NT)}\label{subsub:nt}

The Network Topology signal captures structural anomalies in the
transaction graph that arise from agent enumeration (CHAIN\_1),
disposable agent armies (CHAIN\_4), and cross-platform linkage
(CHAIN\_5).

\paragraph{Sub-signals.}
\emph{Sender centrality} (50\% weight) measures the out-degree and
betweenness centrality of the sending address within a sliding
time window.  An orchestrator spawning 100 disposable agents
(\cref{subsec:chain4}) produces an extreme hub that is
straightforward to detect.
\emph{Receiver centrality} (30\%) applies the same logic to
incoming edges, detecting sink patterns characteristic of value
extraction.
\emph{Path anomaly} (20\%) flags unusual multi-hop paths---for
example, chains of three or more previously unseen addresses
activated within the same block, or repeated edges between the
same pair of addresses.

\paragraph{Relation to attack chains.}
The NT signal is the primary detector for four of the eight chains
in \cref{tab:attack_chains}: CHAIN\_1 (enumeration), CHAIN\_2
(history extraction), CHAIN\_4 (disposable army), and CHAIN\_5
(cross-platform identity).  Its effectiveness stems from a simple
observation: agents must interact with the network to accomplish
their objectives, and those interactions leave structural traces
that are difficult to suppress without fundamentally altering the
attack.


% ---- Signal 3: Value Flow ----
\subsubsection{Signal 3: Value Flow (VF)}\label{subsub:vf}

The Value Flow signal tracks how value moves through the network,
detecting patterns consistent with layering, rapid extraction, and
wash trading.

\paragraph{Sub-signals.}
\emph{Flow velocity} (50\%) measures the rate at which value passes
through an address, normalized by the address's age and transaction
count.  An address that receives and forwards large sums within
minutes scores highly.
\emph{Layering indicator} (50\%) detects fan-out patterns in which
a single inflow is split into multiple outflows of similar size
within a short time window---the canonical layering topology.

\paragraph{A note on net flow imbalance.}
An earlier iteration of the framework included a \emph{net flow
imbalance} sub-signal, hypothesizing that fraudulent agents would
show large imbalances between inflows and outflows.  Phase~5
validation (\cref{sec:evaluation}) revealed a paradoxical result:
the sub-signal fired on 57.1\% of \emph{human} addresses but only
5.6\% of agent addresses.  The explanation is that humans frequently
operate as net senders or net receivers (paying for services,
receiving salary), while agents tend to maintain balanced flows as
part of their operational logic.  This sub-signal was accordingly
removed, illustrating the importance of empirical validation over
intuitive reasoning when designing agent-specific detectors.


% ---- Signal 4: Temporal Consistency ----
\subsubsection{Signal 4: Temporal Consistency (TC)}\label{subsub:tc}

The Temporal Consistency signal exploits the fact that agents
operate on machine time, producing timing distributions that
differ from human behavior in \emph{two} directions: they can be
too fast and too regular.

\paragraph{Sub-signals.}
\emph{Circadian violation} (35\%) flags sustained transaction
activity during hours when human activity is negligible
(typically 02:00--06:00 local time).  While a single late-night
transaction is unremarkable, sustained high-throughput activity
across the circadian trough is a strong agent indicator.
\emph{Timing regularity} (35\%) computes the coefficient of
variation of inter-transaction intervals.  As established in
the CHAIN\_6 analysis (\cref{subsec:chain6}), this sub-signal
fires on \emph{both} tails:
\begin{equation}\label{eq:tc_twotail}
  \text{TC}_{\text{reg}} =
  \begin{cases}
    \text{high} & \text{if } \cvt < 0.05
                  \quad\text{(machine regularity)}, \\
    \text{high} & \text{if inter-tx interval} < 400\;\text{ms}
                  \quad\text{(machine speed)}, \\
    \text{low}  & \text{otherwise}.
  \end{cases}
\end{equation}
The two-tailed design is a direct consequence of the mimicry
paradox: an agent that throttles its timing to appear human
inevitably produces \emph{too little} variance, while an
unthrottled agent produces intervals that are \emph{too short}.
Either direction is diagnostic.
\emph{Burst detection} (30\%) identifies sustained bursts of
activity---defined as $n \geq 10$ transactions with inter-arrival
times below one second---that exceed any plausible human input
rate.


% ---- Signal 5: Cross-Platform Correlation ----
\subsubsection{Signal 5: Cross-Platform Correlation (CP)}\label{subsub:cp}

The Cross-Platform Correlation signal is designed to detect
simultaneous multi-chain presence, the signature of CHAIN\_5.
An entity that transacts on Base, Ethereum mainnet, and Arbitrum
within the same minute---especially if the transactions are
economically linked (same token, complementary amounts)---is
exhibiting behavior that is effectively impossible for a human
operator managing multiple interfaces concurrently.

In the current implementation, CP is \textbf{inactive}: our
validation dataset covers only Base~L2 transactions, and
cross-chain correlation requires oracle infrastructure that is
beyond the scope of this work.  The signal is included in the
architectural specification to ensure the framework generalizes
to the multi-chain setting that CHAIN\_5 anticipates.  Its fusion
weight is accordingly set to zero (\cref{subsec:fusion}).


% ============================================================
\subsection{Signal Fusion}\label{subsec:fusion}

Individual signals capture different facets of agent behavior; the
composite score integrates them into a single risk metric.  We adopt
a linear fusion model with \emph{AUC-proportional} weights---each
signal's weight is proportional to its standalone area under the
ROC curve, measured on the real on-chain validation set of 81,904 USDC
transactions from Base~L2~\cite{dune2026}.

\begin{equation}\label{eq:fusion}
  S_{\text{composite}} =
    0.2739 \cdot \text{NT}
  + 0.2505 \cdot \text{TC}
  + 0.2424 \cdot \text{ER}
  + 0.2332 \cdot \text{VF}
  + 0.0000 \cdot \text{CP}.
\end{equation}

Several properties of \cref{eq:fusion} merit discussion.

\paragraph{Weight interpretation.}
Network Topology carries the largest weight (0.2739), reflecting the
strong structural signatures produced by agent orchestration patterns.
Temporal Consistency is a close second (0.2505), driven by the
two-tailed timing signal that catches both unthrottled and
mimicry-throttled agents.  Economic Rationality and Value Flow
contribute nearly equally (0.2424 and 0.2332, respectively).  The
four active weights sum to approximately one, and the zero weight
on CP ensures that the composite score is unaffected by the
inactive signal.

\paragraph{Effective uniformity of weights.}
We acknowledge that the narrow AUC spread across active signals
(0.4647--0.5990) produces near-uniform weights ($\sigma_w = 0.015$).
AUC-proportional weighting is effectively equivalent to equal
weighting on this dataset: the theoretical justification for
rewarding more discriminative signals has minimal practical effect
when all signals hover near chance.  We considered two alternative
weighting schemes:
\begin{enumerate}[nosep]
  \item \textbf{Exclude TC} (AUC = 0.4647, below chance).  Removing
    TC and re-normalizing yields $w_{\text{NT}} = 0.363$,
    $w_{\text{ER}} = 0.312$, $w_{\text{VF}} = 0.325$---marginally
    less uniform but with a negligible change in composite AUC
    ($+0.003$).  We retain TC because the canary-signal hypothesis
    (\cref{sec:limit_tc_canary}) makes its longitudinal trajectory
    diagnostically valuable even when its current AUC is low.
  \item \textbf{Use $\auc - 0.5$ as weights} (rewarding only
    discriminative power above chance).  This yields
    $w_{\text{NT}} = 0.547$, $w_{\text{VF}} = 0.122$,
    $w_{\text{ER}} = 0.084$, $w_{\text{TC}} = 0.000$
    (below-chance TC is zeroed out), and $w_{\text{CP}} = 0.000$.
    In this regime NT dominates, which is defensible given its
    highest AUC but makes the composite fragile to NT-specific
    evasion.
\end{enumerate}
We retain raw AUC-proportional weighting as a conservative default
that avoids over-concentration on any single signal.  As the
evaluation set grows and AUC values diverge, we expect the weighting
scheme to differentiate naturally; the choice of scheme will matter
more then.

\paragraph{Why linear fusion.}
We deliberately chose a linear model over more expressive alternatives
(gradient-boosted trees, neural networks) for three reasons.
First, interpretability: a linear model allows practitioners to
decompose any flagged transaction into its contributing signals,
which is essential for regulatory compliance and audit trails.
Second, robustness: linear models are less susceptible to
overfitting on the relatively small labeled dataset available for
A2A fraud.
Third, transparency: the weight vector can be published, audited,
and updated independently of the scoring engine, facilitating
community review.

\paragraph{Decision tiers.}
The composite score maps to four operational tiers:

\begin{table}[h]
  \centering
  \caption{Decision tiers derived from composite score thresholds.}
  \label{tab:decision_tiers}
  \small
  \begin{tabular}{@{}l l l@{}}
    \toprule
    \textbf{Tier} & \textbf{Score Range} & \textbf{Action} \\
    \midrule
    \textsc{allow}       & $< 0.09$        & Transaction proceeds normally. \\
    \textsc{flag}        & $0.09$--$0.50$  & Transaction logged for review queue. \\
    \textsc{investigate}  & $0.50$--$0.75$  & Manual or automated investigation triggered. \\
    \textsc{block}       & $\geq 0.75$     & Transaction held pending review. \\
    \bottomrule
  \end{tabular}
\end{table}

\noindent
Threshold selection balances false-positive cost (legitimate
transactions delayed) against false-negative cost (fraudulent
transactions approved).  The \textsc{allow} threshold of 0.09
was chosen to keep the false-positive rate on known-human
traffic below 5\%; the \textsc{block} threshold of 0.75
corresponds to a true-positive rate exceeding 90\% on known-agent
attack traffic in Phase~5 validation (\cref{sec:evaluation}).


% ============================================================
\subsection{Implementation}\label{subsec:implementation}

The detection framework is released as a pip-installable Python
package under the Apache~2.0 license:

\begin{center}
  \texttt{pip install a2a-fincrime-detection}
\end{center}

\paragraph{Architecture.}
The package comprises six signal modules (one per signal, plus a
fusion module), each implementing the sub-signal specifications
of~\cref{subsec:signals}.  All modules follow a common interface:
they accept a transaction DataFrame and return a per-address score
in $[0, 1]$.  The fusion module combines individual scores via
\cref{eq:fusion} and assigns decision tiers per
\cref{tab:decision_tiers}.

\paragraph{Testing and validation.}
The package includes 32 unit tests covering signal computation,
edge cases (single-transaction addresses, zero-value transfers),
and fusion arithmetic.  End-to-end validation was performed on
81,904 on-chain USDC transactions from Base~L2, sourced via Dune
Analytics~\cite{dune2026} and labeled using ERC-8004 agent
registry data~\cite{erc8004}.

\paragraph{Requirements.}
The implementation targets Python $\geq 3.10$ and depends on four
core libraries: \texttt{pandas} for data manipulation,
\texttt{numpy} for numerical computation, \texttt{scikit-learn}
for threshold calibration and ROC analysis, and \texttt{web3} for
on-chain data retrieval.  No GPU is required; scoring the full transaction dataset
completes in under 30 seconds on a single CPU core.

\paragraph{Extensibility.}
Adding a new signal---for example, activating the Cross-Platform
Correlation module when multi-chain data becomes available---requires
implementing the \texttt{SignalModule} interface and adding the
signal's AUC-proportional weight to the fusion vector.  The
architecture is deliberately modular to support the rapid iteration
cycle that A2A fraud detection will demand as the agent ecosystem
evolves.
